\section{Conclusioni}

Tirando le somme, il progetto tocca tutte e quattro le parti richieste: i requisiti, il
modello Asmeta (con gli scenari, sia quelli corretti che quelli nuovi, e il model
checking), i contratti JML sul nucleo del dominio e l'implementazione Java con i vari
test (JUnit, MC/DC, Mockito, model based testing e combinatorial testing), la UI Vaadin,
la CI e l'analisi statica.

\subsection{Riepilogo delle evidenze}

\begin{itemize}
    \item Simulatore e animatore Asmeta (Sezione~2)
    \item AsmetaV -- copertura degli scenari Avalla (Sezione~2)
    \item AsmetaSMV -- 8 CTLSPEC verificate e 2 volutamente false con controesempio (Sezione~2)
    \item Model Advisor e ATGT, opzionali (Sezione~2)
    \item OpenJML ESC sui contratti del nucleo (Sezione~3)
    \item UI Vaadin in esecuzione e test Selenium (Sezione~4)
    \item JUnit (59 test del package \texttt{core}, tutti verdi, + 1 test Selenium UI escluso dalla build standard) e copertura JaCoCo (Sezione~5)
    \item Analisi statica e ispezione critica con LLM (Sezione~6)
    \item GitHub Actions: build e test automatici ad ogni push (Sezione~7)
\end{itemize}

Quello che mi piace di questo percorso è che si tiene tutto insieme: si parte dal
modello astratto, lo si valida e lo si verifica, poi si passa ai contratti JML dimostrati
staticamente e si finisce con i test sul codice e la CI. Alla fine si vede come le
tecniche viste a lezione, ognuna con il suo scopo, entrino a far parte di un unico
processo di sviluppo in cui a ogni livello si controlla qualcosa.
